\section[Cote 121 : la décomposition cellulaire admissible la plus grossière]{Cote 121 : la décomposition cellulaire admissible la plus grossière\protect\verifie{Folder121}}\label{sec:121}

La cote 121, \emph{Topologie modérée}, cherche dans sa première mouture une notion de partie \q{modérée} d'un espace topologique $X$. Les pages~3 à~5 posent une algèbre de Boole $\mathcal F$ de parties de $X$ et les conditions c~bis) ($A \in \mathcal F \Rightarrow \overline A \in \mathcal F$), d) (tout $A \in \mathcal F$ contient un ouvert dense de son adhérence) et e) (toute suite décroissante de fermés de $\mathcal F$, chacun rare dans le précédent, est stationnaire à $\varnothing$) (\lot{121}{1}{4}, \lot{121}{1}{5}). La page~8 montre que la frontière d'un membre est rare (\lot{121}{1}{8}). La page~10 appelle décomposition cellulaire une partition finie de $X$ telle que (1) l'adhérence de chaque cellule soit une réunion de cellules et (2) deux cellules distinctes aient des adhérences distinctes ; elle est $\mathcal F$-admissible si ses cellules sont dans $\mathcal F$ (\lot{121}{1}{10}). La page~14 énonce la proposition suivante et en démontre l'existence par récurrence, en retirant la réunion des frontières et en s'arrêtant par e) (\lot{121}{1}{14}). Comme dans la lecture, on prend dans d) l'intérieur relatif à l'adhérence.

\begin{encadre}
\begin{theoreme}[Grothendieck, page 14\piste{121-coarsest-cell-decomposition}]
Soit $\mathcal F$ une algèbre de Boole de parties d'un espace topologique, stable par adhérence, dont chaque membre contient un ouvert dense de son adhérence, et dans laquelle toute suite décroissante de membres fermés, chacun rare dans le précédent, aboutit à $\varnothing$. Toute sous-famille finie $\mathcal L \subset \mathcal F$ admet une décomposition cellulaire $\mathcal F$-admissible dont chaque élément de $\mathcal L$ est réunion de cellules, et parmi ces décompositions une plus grossière que toutes les autres.
\end{theoreme}
\end{encadre}

\noindent L'énoncé tient : voir \texttt{page14}, avec le lemme de la page~8 (\texttt{Tame.\allowbreak{}rare\_bd}), dans \href{\leanblob/Folder121.lean}{\texttt{Folder121.\allowbreak{}lean}}. La récurrence porte sur un fermé $Z \in \mathcal F$ quelconque à la place de $X$ (\texttt{exists\_coarsest}), la frontière d'un membre étant prise relativement à $Z$. La page affirme sans argument que la décomposition obtenue est la plus grossière ; l'argument de l'édition (\texttt{subset\_bd}) est qu'une cellule d'une décomposition dont $Y$ est réunion de cellules, si elle rencontre la frontière de $Y$, y est contenue, de sorte que toute décomposition admissible convenable se restreint à la partie retirée et raffine les cellules ouvertes.

Ne sont pas démontrés ici : la version à cellules connexes de la page~16, sous la condition f), les corollaires de la page~15 et les décompositions régulières de la page~17.
